\section{BOCPD Configuration and Falsification Detail}
\label{app:bocpd}

This appendix gives the full configuration behind the changepoint hypothesis test of Section \ref{sec:bocpd}. We further stress-test the falsification in three ways, with all intermediate tables released in the accompanying repository. First, replacing the t-approximate $p$ with exact permutation $p$ values removes the apparent significance of the headline cell ($\rho=-0.67$: exact $p=0.080$; Holm-adjusted $0.399$ over the model family), so we report the falsification as pattern evidence---a consistent negative direction across hazards that flips sign at hazard $1/1000$---rather than as a significance claim. Second, removing calendar seasonality before detection (two variants) does not reverse the counterexample: under phase removal traffic keeps the shortest expected run length alongside the longest optimal lookback across all hazards, and under seasonal differencing its rank weakens to third--fourth shortest while the correlation stays non-positive---so seasonality is not the confound behind it. Third, a changepoint-model-free drift measure (rolling-window mean differences and symmetrized Gaussian KL) recovers the negative drift--span relation independently (iTransformer: $\rho \in [-0.90, -0.79]$, strongest $\rho=-0.902$, Holm-adjusted $p=0.033$).

\subsection{Observation model and inference}

We run Bayesian online changepoint detection \citep{bocpd} per dataset with a Gaussian observation model with \emph{unknown mean and unknown variance}: a Normal--Inverse--Gamma prior ($\kappa_0 = 1.0$) conjugate to the likelihood gives a Student-$t$ predictive at every run length. The run-length posterior is truncated at $r_{\max} = 3000$. The hazard is constant, swept over
\begin{equation}
H \in \{1/50,\; 1/100,\; 1/250,\; 1/500,\; 1/1000\}.
\end{equation}
Inference uses the train split only (borders identical to \texttt{data\_provider}), with up to 24 channels subsampled per dataset, and the reported statistic is the time-averaged posterior expected run length $\mathbb{E}[r]$ after a burn-in of 500 steps.

\paragraph{Why the statistic is $\mathbb{E}[r]$, and a discarded degenerate variant.} Under a constant hazard, the per-step posterior changepoint probability $p(r_t{=}0)$ is identically $H$ by construction (both posterior branches scale with the same predictive mass), so it carries no information about the data; the informative, data-driven quantity is the run-length \emph{distribution}, summarized by $\mathbb{E}[r]$, whose downward resets mark detected changepoints. Our first implementation used the common simplification ``unknown mean, known variance''; under a constant hazard that variant is degenerate in the same sense ($p(r_t{=}0) \equiv H$ is a construction identity) and its $\mathbb{E}[r]$ is almost entirely prior-determined, so the likelihood has no influence. We discarded it and report only the NIG (unknown mean \emph{and} variance) model.

\paragraph{Sanity checks.} Two synthetic controls validate the implementation: (i) i.i.d.\ Gaussian input yields $\mathbb{E}[r] \approx 1740$ and keeps growing toward the truncation $r_{\max}=3000$ (no changepoints detected, as expected); (ii) a piecewise mean-shift input yields $\mathbb{E}[r] \approx 102$ with resets at the injected shifts (changepoints detected).

\subsection{Full $\mathbb{E}[r]$ table and rank correlations}

\begin{table}[t]
\caption{Expected run length $\mathbb{E}[r]$ (time-averaged posterior mean, in steps; ETTm* steps are 15 minutes, others hourly/daily) per hazard level $H$, with the measured best lookback under the equal-data-budget protocol (iTransformer / DLinear) for reference. Electricity and traffic --- the shortest $\mathbb{E}[r]$ at moderate $H$ --- have the \emph{longest} optimal lookbacks, the opposite of the run-length hypothesis.}
\label{tab:bocpd_er}
\centering\small
\begin{tabular}{lccccc cc}
\toprule
Dataset & $H{=}1/50$ & $H{=}1/100$ & $H{=}1/250$ & $H{=}1/500$ & $H{=}1/1000$ & Best $L$ (iTr) & Best $L$ (DLin) \\
\midrule
ETTh1 & 65 & 76 & 90 & 102 & 114 & 720 & 1440 \\
ETTh2 & 234 & 248 & 262 & 271 & 279 & 96 & 720 \\
ETTm1 & 60 & 64 & 68 & 71 & 74 & 336 & 336 \\
ETTm2 & 363 & 377 & 390 & 399 & 407 & 336 & 2880 \\
Exchange & 210 & 215 & 219 & 221 & 223 & 96 & 96 \\
Weather & 291 & 297 & 301 & 303 & 305 & 336 & 2880 \\
Electricity & 53 & 68 & 122 & 234 & 372 & 2880 & 2880 \\
Traffic & 48 & 59 & 87 & 146 & 238 & 2880 & 2880 \\
\bottomrule
\end{tabular}
\end{table}

\begin{table}[t]
\caption{Spearman rank correlation between $\mathbb{E}[r]$ and measured best lookback across the eight datasets ($n{=}8$) per hazard level. No hazard yields a significant negative correlation after exact-permutation and multiple-comparison control, and the iTransformer side is negative at $H{=}1/50$ only under a t-approximation (exact permutation $p=0.080$), so we read the table as pattern evidence: the direction is negative at moderate hazards and flips positive at the most extreme hazard on both families (DLinear $+0.723$, exact $p=0.056$).}
\label{tab:bocpd_spearman}
\centering\small
\begin{tabular}{lccc}
\toprule
$H$ & iTransformer & DLinear & Mean of the two \\
\midrule
$1/50$ & $-0.667$ ($p{\approx}0.028$) & $+0.000$ ($p{\approx}1.000$) & $-0.289$ ($p{\approx}0.459$) \\
$1/100$ & $-0.581$ ($p{\approx}0.081$) & $+0.114$ ($p{\approx}0.778$) & $-0.157$ ($p{\approx}0.698$) \\
$1/250$ & $-0.457$ ($p{\approx}0.208$) & $+0.292$ ($p{\approx}0.455$) & $+0.060$ ($p{\approx}0.882$) \\
$1/500$ & $-0.272$ ($p{\approx}0.489$) & $+0.495$ ($p{\approx}0.163$) & $+0.301$ ($p{\approx}0.439$) \\
$1/1000$ & $+0.111$ ($p{\approx}0.784$) & $+0.723$ ($p{\approx}0.010$) & $+0.615$ ($p{\approx}0.056$) \\
\bottomrule
\end{tabular}
\end{table}

Table \ref{tab:bocpd_er} lists $\mathbb{E}[r]$ for every hazard; Table \ref{tab:bocpd_spearman} the rank correlations. The hypothesis ``effective span tracks posterior expected run length'' fails at every hazard level: the two datasets with the shortest $\mathbb{E}[r]$ at moderate hazards (traffic 87, electricity 122 at $H{=}1/250$) are exactly those whose optimal lookback is longest (2880), while ETTm2 has the longest $\mathbb{E}[r]$ (390) but an optimal lookback of only 336. Regime persistence is therefore not the axis that sets the effective span; the long-context benefit on electricity and traffic comes from strong daily/weekly periodicity that persists across regimes, not from regime duration (Section \ref{sec:why}).
